\chapter{Advection-Diffusion-Reaction Solver (Redesign)}
\label{chap:adrRedesign}

\section{Synopsis}
The ADRSolverRedesign solves the same linear advection-diffusion-reaction
problems as the ADRSolver of the previous chapter -- the governing equations
are unchanged -- on the redesigned operator framework, in which every step of
the discretisation (basis transforms, volume and trace fluxes, boundary
conditions, the linear solve) is an operator with interchangeable
implementations for each execution space. The same session file therefore runs

\begin{itemize}
\item in \inltt{Serial} on a single core,
\item in \inltt{AVX} using SIMD vectorisation (AVX2/AVX512 on x86, NEON on
  Arm), and
\item on a \inltt{Device} (CUDA, HIP or SYCL GPU back-ends),
\end{itemize}

and in parallel under MPI, where the discontinuous Galerkin trace exchange
overlaps interior work with communication.

The solver currently provides:

\begin{itemize}
\item the unsteady advection, diffusion and advection-diffusion problems
  and the steady advection-diffusion-reaction, Helmholtz and Poisson
  problems of Table~\ref{t:adrredesign:eqtype}, in 1D, 2D and 3D on every
  element shape;
\item both a continuous (\inltt{Continuous}) and a discontinuous
  (\inltt{DisContinuous}) Galerkin projection;
\item for DG, weak-DG advection with an upwind trace flux and interior-penalty
  diffusion, integrated explicitly;
\item for CG, non-conservative advection and an implicit treatment of
  diffusion through the iterative linear solvers of the framework;
\item an anisotropic diffusion tensor and a reaction coefficient;
\item time-dependent Dirichlet boundary conditions, and restart from
  checkpoint (\inltt{.fld}) files.
\end{itemize}

Features of the legacy solver not yet available here include LDG diffusion,
implicit or IMEX time integration of a DG discretisation, the Burgers and
moving-frame (\inltt{MMF}) systems, and the \inltt{Projection} system.

\section{Usage}
\begin{lstlisting}[style=BashInputStyle]
ADRSolverRedesign session.xml
\end{lstlisting}

The execution space may be chosen on the command line without touching the
session file:

\begin{lstlisting}[style=BashInputStyle]
ADRSolverRedesign --opExecSpace AVX session.xml
\end{lstlisting}

The implementation used within a space (\inltt{SumFac} or \inltt{StdMat}
element kernels) can likewise be forced for every operator at once with
\inltt{--opImpl}. When no space is given, the solver picks the most capable one
the build supports and says which. Individual operators may instead be given
their own implementation from the session file, in an \inltt{<OPERATORS>}
block with one section per back-end whose children name an operator and
contain its implementation:

\begin{lstlisting}[style=XmlStyle]
<OPERATORS>
  <AVXBACKEND>
    <BWDTRANS> SumFac </BWDTRANS>
    <PHYSDERIV> StdMat </PHYSDERIV>
  </AVXBACKEND>
</OPERATORS>
\end{lstlisting}

\section{Session file configuration}
\label{sec:sessionFileADRRedesign}
The session file follows the structure described for the legacy solver; below
we describe the options currently understood, and where they differ.

\subsection*{Solver info}
\begin{lstlisting}[style=XmlStyle]
<SOLVERINFO>
  <I PROPERTY="EQTYPE"                VALUE="UnsteadyAdvectionDiffusion" />
  <I PROPERTY="Projection"            VALUE="DisContinuous"              />
  <I PROPERTY="AdvectionType"         VALUE="WeakDG"                     />
  <I PROPERTY="UpwindType"            VALUE="Upwind"                     />
  <I PROPERTY="DiffusionType"         VALUE="InteriorPenalty"            />
  <I PROPERTY="AdvectionAdvancement"  VALUE="Explicit"                   />
  <I PROPERTY="DiffusionAdvancement"  VALUE="Explicit"                   />
</SOLVERINFO>
\end{lstlisting}
\begin{itemize}
\item \inltt{EQTYPE} selects the system from
  Table~\ref{t:adrredesign:eqtype}. A reaction term enters any of them
  through the parameter \inltt{lambda}.
\item \inltt{Projection} is \inltt{Continuous} or \inltt{DisContinuous}. The
  two projections have different capabilities, set out under the two
  advancement properties below.
\item \inltt{AdvectionType} is \inltt{WeakDG} for a discontinuous projection
  and \inltt{NonConservative} for a continuous one.
\item \inltt{UpwindType} names the trace flux of the weak-DG advection;
  \inltt{Upwind} is the one available.
\item \inltt{DiffusionType} is \inltt{InteriorPenalty}, the scheme the
  discontinuous projection uses for the diffusion term; a continuous
  projection needs no such choice.
\item \inltt{AdvectionAdvancement} must be \inltt{Explicit}: advection is not
  integrated implicitly.
\item \inltt{DiffusionAdvancement} is \inltt{Explicit} or \inltt{Implicit}.
  With a discontinuous projection it must be \inltt{Explicit}; with a
  continuous one \inltt{Implicit} is available and solves the diffusion
  operator with the iterative linear solvers below, under an IMEX or
  diagonally implicit time integration scheme.
\item \inltt{GLOBALSYSSOLN}, \inltt{LinSysIterSolver} and
  \inltt{Preconditioner} configure the linear solve of a continuous
  projection, for implicit diffusion and for the steady systems:
  \inltt{IterativeFull}; \inltt{ConjugateGradientLoc} for a symmetric system
  or \inltt{GMRES} where advection makes it non-symmetric (the solver switches
  to GMRES itself, with a warning, if asked for conjugate gradients on an
  advection system); and \inltt{Diagonal}.
\end{itemize}

\begin{table}[h]
\begin{center}
\begin{tabular}{lll}
\toprule
\inltt{EQTYPE} & Problem & Time dependence \\
\midrule
\inltt{UnsteadyAdvection} & advection & unsteady \\
\inltt{UnsteadyDiffusion} & diffusion & unsteady \\
\inltt{UnsteadyAdvectionDiffusion} & advection-diffusion & unsteady \\
\inltt{SteadyADR} & advection-diffusion-reaction & steady \\
\inltt{Helmholtz} & Helmholtz & steady \\
\inltt{Poisson} & Poisson & steady \\
\bottomrule
\end{tabular}
\end{center}
\caption{Systems solved by the redesigned ADR solver.}
\label{t:adrredesign:eqtype}
\end{table}

\subsection*{Parameters}
\begin{lstlisting}[style=XmlStyle]
<PARAMETERS>
  <P> TimeStep                 = 0.001             </P>
  <P> NumSteps                 = 100               </P>
  <P> FinTime                  = NumSteps*TimeStep </P>
  <P> IO_InfoSteps             = 10                </P>
  <P> epsilon                  = 0.01              </P>
  <P> lambda                   = 0.0               </P>
  <P> IterativeSolverTolerance = 1e-12             </P>
</PARAMETERS>
\end{lstlisting}
\begin{itemize}
\item \inltt{TimeStep}, \inltt{NumSteps}, \inltt{FinTime} and
  \inltt{IO\_InfoSteps} behave as in the legacy solver. The step is fixed.
\item \inltt{epsilon} sets the diffusion coefficient $\epsilon$. Default 1.
\item \inltt{D00}, \inltt{D01}, \inltt{D11}, \inltt{D02}, \inltt{D12},
  \inltt{D22} set the entries of a symmetric diffusion tensor $D$, as many as
  the dimension needs; when any is given the tensor replaces $\epsilon$.
  Default: the identity.
\item \inltt{lambda} sets the reaction coefficient $\lambda$. Default 0.
\item \inltt{IPPenaltyCoeff} multiplies the interior-penalty jump term. With
  its value $c$, the penalty across a trace is $c\,\max(p+1)^2/h$, with $p+1$
  the larger number of modes of the two elements sharing it and $h$ the
  element length normal to it. Default 1, which is the penalty the legacy
  solver applies.
\item \inltt{IPSymmFluxCoeff} and \inltt{IP2ndDervCoeff} select the variant
  of the interior-penalty scheme. Both default to 0, the incomplete scheme.
\item \inltt{IterativeSolverTolerance} is the residual tolerance of the
  iterative linear solve of a continuous projection.
\end{itemize}

\subsection*{Functions}
\begin{itemize}
\item \inltt{AdvectionVelocity} specifies the advection velocity
  $\mathbf{V}$, one component per coordinate direction.
\item \inltt{InitialConditions} specifies the initial condition of an
  unsteady problem, as an expression or as a checkpoint file to restart from.
\item \inltt{BodyForce} specifies the forcing function $f$. The legacy solver
  calls this function \inltt{Forcing}.
\item \inltt{ExactSolution} is used to report the error at the end of the run.
\end{itemize}

\subsection*{Time integration scheme}
The \inltt{<TIMEINTEGRATIONSCHEME>} block is unchanged from the legacy solver.
A discontinuous projection needs an explicit family (\inltt{RungeKutta},
including its SSP variants, \inltt{ExplicitSDC}); a continuous projection with
implicit diffusion takes \inltt{IMEX} when advection is present and a
diagonally implicit scheme such as \inltt{BDFImplicit} when it is not.

\section{Examples}
The regression tests under \inltt{solvers/ADRSolverRedesign/Tests} are
complete, runnable session files; two are quoted here.

\subsection*{Weak-DG advection with a time-dependent Dirichlet condition}
From \inltt{UnsteadyAdvection1D\_TimeDependentDirichlet\_DG.xml}: a 1D
advection problem, integrated with fourth-order Runge-Kutta, whose boundary
value is a function of time.
\begin{lstlisting}[style=XmlStyle]
<SOLVERINFO>
  <I PROPERTY="EQTYPE"               VALUE="UnsteadyAdvection" />
  <I PROPERTY="Projection"           VALUE="DisContinuous"     />
  <I PROPERTY="AdvectionType"        VALUE="WeakDG"            />
  <I PROPERTY="UpwindType"           VALUE="Upwind"            />
  <I PROPERTY="AdvectionAdvancement" VALUE="Explicit"          />
</SOLVERINFO>
<TIMEINTEGRATIONSCHEME>
  <METHOD> RungeKutta </METHOD>
  <ORDER> 4 </ORDER>
</TIMEINTEGRATIONSCHEME>
<PARAMETERS>
  <P> TimeStep       = 0.001              </P>
  <P> NumSteps       = 100                </P>
  <P> FinTime        = NumSteps*TimeStep  </P>
  <P> advx           = 1.0                </P>
</PARAMETERS>
\end{lstlisting}

\subsection*{Interior-penalty diffusion on a hybrid mesh}
From \inltt{ExDiffusion\_2D\_IP\_hybrid\_m3.xml}: explicit diffusion on a mesh
of triangles and quadrilaterals, with the interior-penalty scheme at its
default coefficients.
\begin{lstlisting}[style=XmlStyle]
<SOLVERINFO>
  <I PROPERTY="EQTYPE"                VALUE="UnsteadyDiffusion" />
  <I PROPERTY="Projection"            VALUE="DisContinuous"     />
  <I PROPERTY="AdvectionType"         VALUE="WeakDG"            />
  <I PROPERTY="DiffusionType"         VALUE="InteriorPenalty"   />
  <I PROPERTY="DiffusionAdvancement"  VALUE="Explicit"          />
</SOLVERINFO>
<TIMEINTEGRATIONSCHEME>
  <METHOD> RungeKutta </METHOD>
  <ORDER> 4 </ORDER>
</TIMEINTEGRATIONSCHEME>
<PARAMETERS>
  <P> TimeStep        = 0.000001          </P>
  <P> NumSteps        = 500               </P>
  <P> FinTime         = TimeStep*NumSteps </P>
</PARAMETERS>
\end{lstlisting}
The companion files \inltt{\_Neumann}, \inltt{\_mixedBC} and
\inltt{\_periodic} run the same problem under those boundary conditions, and
\inltt{ExDiffusion\_3D\_IP\_prism\_hex.xml} the 3D case on prisms and
hexahedra.
